\section{Why the maximum occurs at a corner}
\label{sec:argmax}

The location of the maximum is governed by two different effects.
Translation of the angular mesh produces a common upper envelope, while
reflection of the interpolation kernel produces cancellation.  The
envelope is largest near the odd vertices of the angular mesh.  At an
interior odd vertex, however, some reflected terms have opposite signs;
this loss excludes the vertex and its neighborhood.  At the two odd
corners the reflected terms coincide.  There we compare the Lebesgue
function itself with its corner value.

We state the result with coordinates, since different orientations of the
Padua points occur in the literature.  Throughout this section,
\begin{equation}
 \begin{split}
 \mathcal P_n=\Bigl\{\left(\cos\frac{j\pi}{n},
                  \cos\frac{k\pi}{n+1}\right):&\ 0\leq j\leq n,\\
 &0\leq k\leq n+1,\ j+k\equiv0\pmod2\Bigr\},\quad n\geq2.
 \end{split}                                         \label{eq:argmax-nodes}
\end{equation}
Let $L_{jk}\in\Pi_n$ be its cardinal polynomials, where $\Pi_n$ denotes
the bivariate polynomials of total degree at most $n$, and put
\[
 \lambda_n(z)=\sum_{j+k\ \text{even}}|L_{jk}(z)|,
 \qquad \Lambda_c(n)=\lambda_n(-1,-1).
\]

\begin{theorem}[Exact corner maximizers]
\label{thm:argmax-all-degrees}
For every integer $n\geq2$,
\[
 \max_{z\in[-1,1]^2}\lambda_n(z)=\Lambda_c(n),
 \qquad
 \mathop{\rm argmax}_{[-1,1]^2}\lambda_n=
 \begin{cases}
 \{(-1,-1),(1,-1)\},&n\text{ even},\\
 \{(-1,-1),(-1,1)\},&n\text{ odd}.
 \end{cases}                                         \label{eq:argmax-exact-set}
\]
The assertion is computer assisted: the proof for $n\geq320$ uses the
quantitative limiting-constant enclosure specified in
Appendix~\ref{app:argmax-details}; degrees $2\leq n\leq319$ are covered
by the exact certificates described in Section~\ref{sec:certificates}.
\end{theorem}

Here and below, the corner value in a comparison is the exact value
$\Lambda_c(n)$.  A rational lower endpoint retained by a checker is used
only to establish a strict inequality against that value.

\subsection{A common envelope and the cancellation it omits}

Put $m=n+1$, $\mathsf N=nm$, $h_0=1$, and $h_r=2$ for $r>0$.  Define the
finite trigonometric polynomial
\begin{equation}
 G_n(u,v)=\sum_{r+s\leq n}h_rh_s\cos(ru)\cos(sv)-\cos(nu).
                                                        \label{eq:argmax-kernel}
\end{equation}
In particular, the coefficient of its exceptional mode $(n,0)$ is one.
Finite sums define all values, including collisions in alternative
quotient formulas for the kernel.  Its common envelope is
\begin{equation}
 S_n(\alpha,\beta)=\frac1{2\mathsf N}
 \sum_{\substack{0\leq r<2n,\ 0\leq s<2m\\r+s\ {\rm odd}}}
 \left|G_n\left(\frac{r\pi+\alpha}{n},
                    \frac{s\pi+\beta}{m}\right)\right|.
                                                        \label{eq:argmax-primitive}
\end{equation}
The angular coordinates of a physical point are written
\begin{equation}
 \theta=\frac{J\pi+\alpha}{n},\qquad
 \phi=\frac{K\pi+\beta}{m},\qquad
 J+K\ \text{odd},\qquad \rho=|\alpha|+|\beta|\leq\pi,
                                                        \label{eq:argmax-chart}
\end{equation}
with $0\leq\theta,\phi\leq\pi$, $0\leq J\leq n$, and $0\leq K\leq m$.
The indices $J,K$ locate an odd mesh vertex; they need not be node
indices.  Several such charts may represent a point on a shared boundary.

\begin{lemma}[Physical cover and exact defect]
\label{lem:argmax-physical-cover}
The charts \eqref{eq:argmax-chart} cover the closed physical square.  In
each admissible chart there is a nonnegative full reflection defect
$\Delta_{JK}(\alpha,\beta)$ such that
\begin{equation}
 \lambda_n(\cos\theta,\cos\phi)
       =S_n(\alpha,\beta)-\Delta_{JK}(\alpha,\beta),
 \qquad \Delta_{JK}\geq\max\{Q_h(J),Q_v(K)\}\geq0.
                                                        \label{eq:argmax-defect}
\end{equation}
The quantities $Q_h,Q_v$ are the losses obtained by pairing reflected
terms in one coordinate.  Each interior index supplies a positive local
lower bound for the corresponding projection.  At a physical nonnode
corner $\Delta=0$, and $S_n(0,0)=\Lambda_c(n)$.
\end{lemma}

\begin{proof}
The cardinal formula expresses $L_{jk}$ as the weighted average of four
reflected copies of $G_n$.  Summing the absolute values before those
four terms are added gives $S_n$; their triangle-inequality loss is
$\Delta_{JK}$.  Half weights at an endpoint exactly compensate the two
coincident reflected preimages.  Pairing in one coordinate gives the
two separate lower bounds in \eqref{eq:argmax-defect}.

For coverage, choose an angular mesh rectangle containing the point.
In coordinates scaled to that rectangle, the distances to its two odd
vertices are either $a+b,2-a-b$ or $1-a+b,1+a-b$.  They sum to two,
so one is at most one.  This gives $\rho\leq\pi$, including all ties.
At a physical corner all four copies of the kernel coincide.  The
explicit identities and normalizations are given in
Appendix~\ref{app:argmax-details}.
\end{proof}

The distinction between $S_n$ and $\lambda_n$ is essential.  The common
envelope can increase initially in one phase direction.  The defect
excludes this increase at a noncorner center.  At a corner we use a
direct local descent estimate instead of assigning a nonexistent
positive defect there.

The delicate contribution comes from two channels of nearly coincident
kernel poles.  Moving the phases changes which of two opposing terms is
smaller, so taking separate upper bounds too early loses useful
cancellation.  We first integrate the complete paired contribution and
then bound it.  In the first phase direction the surviving upper term
has order $x^{3/2}$, rather than quadratic order, and can dominate a
quadratic envelope loss close to zero.  The reflection defect pays for
this increase at noncorner centers.  At an actual corner we instead
prove descent directly for the physical cardinal sum.

\subsection{Four overlapping regions suffice}

The analytic estimates have a simple geometric organization.  Write
$x=|\alpha|$, $y=|\beta|$, and $\rho=x+y$.

\begin{lemma}[Quantitative comparisons]
\label{lem:argmax-four-regions}
For every integer $n\geq320$ the following statements hold on their
entire closed domains:
\begin{align}
 \lambda_n&<\Lambda_c(n)-\frac1{20000},
 &&\rho\leq\tfrac12,\quad (J,K)\text{ noncorner};
                                                        \label{eq:argmax-half-near}\\
 \lambda_n&<\Lambda_c(n)-\frac3{1000},
 &&\tfrac12\leq\rho\leq\tfrac35,\quad (J,K)\text{ noncorner};
                                                        \label{eq:argmax-annulus}\\
 S_n&<\Lambda_c(n)-\frac{99}{2500},
 &&\tfrac35\leq\rho\leq1;
                                                        \label{eq:argmax-wide}\\
 S_n&<\Lambda_c(n)-\frac{13}{500},
 &&1\leq\rho\leq\pi.                               \label{eq:argmax-outer}
\end{align}
At either parity-correct nonnode corner, in inward coordinates,
\begin{equation}
 \lambda_n\left(-\cos\frac A n,-\cos\frac Bm\right)
 \leq\Lambda_c(n)-\frac{A^2+B^2}{30},
 \qquad 0\leq A,B\leq\frac\pi4.                    \label{eq:argmax-corner-loss}
\end{equation}
The last assertion already holds for every $n\geq256$, with the reflected
coordinate chart at the other corner.
\end{lemma}

The proof of this lemma, including its explicit constants, is in
Appendix~\ref{app:argmax-details}.  Near the origin it combines a
fractional-power upper bound for the central kernel contribution with
a positive reflection defect.  On the short annulus a smaller defect
still suffices.  Beyond $\rho=3/5$ the common envelope alone is strictly
below the corner value.  The final corner estimate comes from
interpolating the smooth signed cardinal sums in squared inward
coordinates before taking their absolute values.  This avoids assuming
smoothness of the Lebesgue function across a cardinal zero.

\begin{theorem}[Analytic tail]
\label{thm:argmax-analytic-tail}
Equation~\eqref{eq:argmax-exact-set} holds for every integer $n\geq320$.
\end{theorem}

\begin{proof}
Represent any physical point by Lemma~\ref{lem:argmax-physical-cover}.
At a noncorner center, \eqref{eq:argmax-half-near} and
\eqref{eq:argmax-annulus} exclude $\rho\leq3/5$.  At larger radii,
\eqref{eq:argmax-defect}, \eqref{eq:argmax-wide}, and
\eqref{eq:argmax-outer} give strict exclusion at every center.  If the
center is an actual corner and $\rho\leq3/5$, its admissible phases are
inward and belong to the square in \eqref{eq:argmax-corner-loss}, since
$3/5<\pi/4$.  Equality is possible there only at $A=B=0$.

The endpoint pairs with odd index sum are exactly $(n,m),(0,m)$ for
even $n$, and $(n,m),(n,0)$ for odd $n$.  They give precisely the two
corners displayed in \eqref{eq:argmax-exact-set}.  For even $n$ the
reflection $x\mapsto-x$ permutes the nodes and their cardinal
polynomials; for odd $n$, $m$ is even and $y\mapsto-y$ does so.
Consequently their values agree exactly.  They attain the maximum,
which also exists by continuity on the compact square.  Closed chart
ties, axes, node values, and endpoints have all been included in the
same exclusions.
\end{proof}

\begin{proof}[Proof of Theorem~\ref{thm:argmax-all-degrees}]
Theorem~\ref{thm:argmax-analytic-tail} covers $n\geq320$.
The finite certificates of Section~\ref{sec:certificates} cover every
integer $2\leq n\leq319$ and prove the identical strict equality
classification.  These two degree ranges exhaust the assertion.
\end{proof}
